\section{Pairs with Cut-Value at Most Two}
\label{sec:squares}

\subsection{Minimum Cuts of a Pair with Cut-Value Two}
\label{sec:structure}

The minimum $u$--$v$ cuts of a pair with $c(u, v) = 2$ arrange the graph into a chain of subgraphs, the squares, in which consecutive squares are joined in two ways (\cref{thm:structure}, \cref{fig:square-decomposition}).
The structure of all minimum cuts between two vertices is known in general form, as the acyclic digraph of Picard and Queyranne~\cite{PicardQueyranne1980} and, for undirected graphs, as the strip of Dinitz and Vainshtein (see~\cite{BaswanaPandey2025Carcass} and \cref{sec:related-work}).
\Cref{lem:orientation,lem:pieces,lem:degrees} and statements (i) to (iv) of \cref{thm:structure} describe this structure for cut-value two in the explicit form that the later sections use, with squares, links, general paths and terminals.
We include their proofs to keep the paper self-contained.
Statement (v) of \cref{thm:structure} corresponds to the construction of the cactus that represents the cuts with two edges of a 2-edge-connected graph, in which every class of the relation $c(x, y) \geq 3$ is contracted to one node~\cite{Mehlhorn2017Certifying}.

Throughout this section and \cref{sec:separating}, we assume that $c(u, v) = 2$.
By Menger's theorem~\cite[Corollary~3.3.5]{Diestel2025}, $G$ contains two edge-disjoint $u$--$v$ paths.
We fix two such paths $p_1$ and $p_2$.
Let $F = F(u, v)$ be the union of all minimum $u$--$v$ cuts.

\begin{lemma}
\label{lem:orientation}
(i) Every minimum $u$--$v$ cut consists of one edge of $p_1$ and one edge of $p_2$, so $F \subseteq p_1 \cup p_2$.
(ii) Let $\delta(A)$ be a minimum $u$--$v$ cut and let $R$ be a $u$--$v$ path.
Then $R$ contains exactly one edge of $\delta(A)$, the vertices of $R$ before this edge lie in $A$, and the vertices of $R$ after it lie outside $A$.
(iii) Every edge $e \in F$ has an end that lies in $A$ for every minimum $u$--$v$ cut $\delta(A)$ that contains $e$, and every $u$--$v$ path that contains $e$ meets this end before the other end of $e$.
\end{lemma}

\begin{proof}
(i) The graph obtained by removing a minimum $u$--$v$ cut contains neither $p_1$ nor $p_2$.
The cut therefore contains an edge of each path, and these two edges are distinct.

(ii) By \cref{lem:minimum-cuts}, the path $R$ contains an odd number of the two edges of $\delta(A)$, hence exactly one.
The path starts in $A$, ends outside $A$ and changes between $A$ and $V \setminus A$ only at this edge.

(iii) By (i), the edge $e$ lies on $p_j$ for some $j \in \{1, 2\}$.
Let $s$ be the end of $e$ that $p_j$ meets first.
For every minimum $u$--$v$ cut $\delta(A)$ that contains $e$, statement (ii) for the path $p_j$ gives $s \in A$, and the other end of $e$ lies outside $A$.
Statement (ii) for a $u$--$v$ path $R$ that contains $e$ and for one such cut shows that $R$ meets $s$ before the other end of $e$.
\end{proof}

\begin{definition}[Pieces, squares and links]
\label{def:pieces}
Let $e \in F$.
The end of $e$ that lies in $A$ for every minimum $u$--$v$ cut $\delta(A)$ containing $e$ is the \emph{tail} of $e$, and the other end is the \emph{head} of $e$.
A \emph{piece} is a component of $G \setminus F$.
The edge $e$ \emph{leaves} the piece that contains its tail and \emph{enters} the piece that contains its head.
A piece is a \emph{square} if it contains $u$, if it contains $v$, or if two edges of $F$ enter it.
Every other piece is a \emph{link}.
\end{definition}

By \cref{lem:orientation}(iii), every $u$--$v$ path that contains an edge of $F$ traverses it from its tail to its head.
The tail of an edge of $F \cap p_j$ is the end that $p_j$ meets first, but the definition does not depend on the choice of $p_1$ and $p_2$.
We use the same symbol for a piece and for its vertex set, and we say that a path \emph{visits} a piece if it contains a vertex of the piece.

\begin{lemma}
\label{lem:pieces}
(i) For every minimum $u$--$v$ cut $\delta(A)$, every piece is contained in $A$ or disjoint from $A$.
(ii) The two ends of every edge of $F$ lie in different pieces, and every edge of $G$ whose ends lie in different pieces belongs to $F$.
(iii) Let $R$ be a $u$--$v$ path that visits a piece $X$.
Then the vertices of $R$ in $X$ are consecutive on $R$.
If $u \notin X$, the edge of $R$ before the first of these vertices enters $X$, and if $v \notin X$, the edge of $R$ after the last of them leaves $X$.
No other edge of $F$ on $R$ has an end in $X$.
\end{lemma}

\begin{proof}
(i) Since $\delta(A) \subseteq F$, every piece is a connected subgraph of $G \setminus \delta(A)$, and no edge of $G \setminus \delta(A)$ joins $A$ to $V \setminus A$.

(ii) An edge of $F$ lies in a minimum $u$--$v$ cut $\delta(A)$ and has one end in $A$ and one outside $A$, so its ends lie in different pieces by (i).
An edge outside $F$ is an edge of $G \setminus F$, and its ends lie in the same component of this graph.

(iii) Let $x$ be a vertex of $R$ in $X$ whose successor $y$ on $R$ is not in $X$.
By (ii), $xy \in F$, and by \cref{lem:orientation}(iii), $x$ is the tail of $xy$, so $xy$ leaves $X$.
Let $\delta(A)$ be a minimum $u$--$v$ cut that contains $xy$.
Then $x \in A$, hence $X \subseteq A$ by (i), and by \cref{lem:orientation}(ii) all vertices of $R$ after $x$ lie outside $A$ and therefore outside $X$.
Thus $x$ is the last vertex of $R$ in $X$.
It follows that the vertices of $R$ in $X$ are consecutive and that, if $v \notin X$, the edge of $R$ after the last of them leaves $X$.
If $u \notin X$, let $x'$ be the first vertex of $R$ in $X$ and let $y'$ be its predecessor.
By (ii), $y'x' \in F$, and by \cref{lem:orientation}(iii), $x'$ is its head, so $y'x'$ enters $X$.
Finally, an edge of $F$ on $R$ with an end in $X$ has its other end outside $X$ by (ii), so it is one of the two edges just described.
\end{proof}

By \cref{lem:pieces}(ii), a piece is an induced subgraph of $G$.
For a $u$--$v$ path $R$ that visits a piece $X$, the \emph{trace} of $R$ in $X$ is the path formed by the vertices of $R$ in $X$, which is a path of $X$ by \cref{lem:pieces}(ii) and~(iii).

\begin{lemma}
\label{lem:degrees}
(i) The vertices $u$ and $v$ lie in different pieces.
(ii) Let $j \in \{1, 2\}$ and let $X$ be a piece.
If $p_j$ visits $X$ and $u \notin X$, exactly one edge of $p_j$ enters $X$, and otherwise none does.
If $p_j$ visits $X$ and $v \notin X$, exactly one edge of $p_j$ leaves $X$, and otherwise none does.
(iii) Every piece is visited by $p_1$ or by $p_2$.
Consequently, no edge enters the piece of $u$ and two edges leave it, and two edges enter the piece of $v$ and none leaves it.
A link is a piece that is visited by exactly one of $p_1$ and $p_2$, and it has exactly one entering and one leaving edge.
Every square is visited by both paths.
\end{lemma}

\begin{proof}
(i) Since $c(u, v) = 2$, there is a minimum $u$--$v$ cut $\delta(A)$, and by \cref{lem:pieces}(i) the piece of $u$ lies in $A$ and the piece of $v$ outside $A$.

(ii) If $p_j$ does not visit $X$, no edge of $p_j$ has an end in $X$.
Otherwise the statement is \cref{lem:pieces}(iii) for $R = p_j$.

(iii) Let $X$ be a piece.
By (i), $X \neq V$, and since $G$ is connected, some edge has exactly one end in $X$.
This edge belongs to $F$ by \cref{lem:pieces}(ii) and lies on $p_1$ or $p_2$ by \cref{lem:orientation}(i), so one of the two paths visits $X$.

The consequences follow from (ii) and (iii).
Every edge of $F$ lies on exactly one of $p_1$ and $p_2$, and both paths visit the pieces of $u$ and of $v$.
A piece that contains neither $u$ nor $v$ is entered by as many edges as there are paths among $p_1, p_2$ that visit it, and as many edges leave it.
\end{proof}

By \cref{lem:degrees}, both paths visit every square.
Every square that does not contain $u$ is entered by exactly one edge of $p_1$ and exactly one edge of $p_2$, and every square that does not contain $v$ is left by exactly one edge of each of the two paths.

\begin{definition}[General paths and terminals]
\label{def:decomposition}
Let $S_1, \dots, S_k$ be the squares in the order in which $p_1$ visits them.
For $2 \leq i \leq k$, the vertices $a_i$ and $b_i$ are the heads of the edges of $p_1$ and of $p_2$ that enter $S_i$, and $a_1 = b_1 = u$.
For $1 \leq i < k$, the vertices $c_i$ and $d_i$ are the tails of the edges of $p_1$ and of $p_2$ that leave $S_i$, and $c_k = d_k = v$.
The vertices $a_i, b_i, c_i, d_i$ are the \emph{terminals} of $S_i$, the first two are its \emph{entering terminals} and the last two its \emph{leaving terminals}.
A \emph{general path} from a square $S$ to a square $S'$ is a sequence $e_0, L_1, e_1, \dots, L_r, e_r$ with $r \geq 0$ in which $L_1, \dots, L_r$ are links, the edge $e_0$ leaves $S$, the edge $e_r$ enters $S'$, and for $1 \leq m \leq r$ the edge $e_{m-1}$ enters $L_m$ and the edge $e_m$ leaves $L_m$.
\end{definition}

We have $u \in S_1$ and $v \in S_k$, and $k \geq 2$ by \cref{lem:degrees}(i).
The four terminals of a square need not be distinct, and by \cref{lem:orientation}(i) each of them lies on $p_1$ or on $p_2$.

\begin{theorem}[Square decomposition]
\label{thm:structure}
Let $G$ be a connected graph, let $u$ and $v$ be two vertices of $G$ with $c(u, v) = 2$, and let $p_1$ and $p_2$ be two edge-disjoint $u$--$v$ paths.
With the notions of \cref{def:pieces,def:decomposition}, the following statements hold.
(i) For $2 \leq i \leq k$, the two edges that enter $S_i$ form a minimum $u$--$v$ cut, and for $1 \leq i < k$, so do the two edges that leave $S_i$.
(ii) Every $u$--$v$ path visits all squares, in the order $S_1, \dots, S_k$.
(iii) For $1 \leq i < k$ and $j \in \{1, 2\}$, there is exactly one general path $P_{i,j}$ from $S_i$ to $S_{i+1}$ whose edges lie on $p_j$, and every link and every edge of $F$ belongs to exactly one of these $2(k - 1)$ general paths.
(iv) For $1 \leq i < k$, every $u$--$v$ path contains all edges and visits all links of one of the general paths $P_{i,1}$ and $P_{i,2}$, and it contains no edge and visits no link of the other one.
(v) Any two distinct terminals $s$ and $t$ of the same square satisfy $c(s, t) \geq 3$.
\end{theorem}

\begin{proof}
(i) Let $2 \leq i \leq k$, and let $g_1$ and $g_2$ be the edges of $p_1$ and $p_2$ that enter $S_i$.
For $j \in \{1, 2\}$, choose a minimum $u$--$v$ cut $\delta(A_j)$ that contains $g_j$.
The head of $g_j$ lies in $S_i$ and outside $A_j$, so $S_i$ is disjoint from $A_j$ by \cref{lem:pieces}(i).
The sets $A_1 \cap A_2$ and $A_1 \cup A_2$ contain $u$ and not $v$, so $\delta(A_1 \cap A_2)$ and $\delta(A_1 \cup A_2)$ are $u$--$v$ cuts.
Each of them has at least two edges, and by \cref{lem:uncrossing} they have at most four edges together.
Hence $\delta(A)$ with $A = A_1 \cup A_2$ is a minimum $u$--$v$ cut.
The tail of $g_j$ lies in $A_j \subseteq A$ and its head lies in $S_i$, which is disjoint from $A$, so $\delta(A) = \{g_1, g_2\}$.
For $1 \leq i < k$ and the edges $h_1, h_2$ that leave $S_i$, let $\delta(B_j)$ be a minimum $u$--$v$ cut that contains $h_j$.
Now $S_i \subseteq B_j$, and the same argument shows that $\delta(B)$ with $B = B_1 \cap B_2$ is a minimum $u$--$v$ cut.
The tail of $h_j$ lies in $S_i \subseteq B$ and its head lies outside $B_j$, so $\delta(B) = \{h_1, h_2\}$.

(ii) Let $R$ be a $u$--$v$ path.
It visits $S_1$ and $S_k$.
Let $2 \leq i \leq k$, and let $\delta(A)$ be the minimum $u$--$v$ cut formed by the edges that enter $S_i$.
By \cref{lem:orientation}(ii), $R$ contains exactly one edge of $\delta(A)$, the vertices of $R$ before it lie in $A$ and the vertices after it outside $A$.
The head of this edge lies in $S_i$, so $R$ visits $S_i$, and this head is the first vertex of $R$ in $S_i$.
Let $X$ be a square other than $S_i$.
By \cref{lem:pieces}(i), $X \subseteq A$ or $X$ is disjoint from $A$.
In the first case $R$ visits $X$ before $S_i$, and in the second case after $S_i$, because the vertices of $R$ in $S_i$ are consecutive.
Which case occurs does not depend on $R$.
Hence all $u$--$v$ paths visit the squares in the same order as $p_1$.

(iii) Let $1 \leq i < k$ and $j \in \{1, 2\}$.
By (ii), the path $p_j$ visits $S_i$, then $S_{i+1}$, and no square in between.
Let $X_0 = S_i, X_1, \dots, X_r, X_{r+1} = S_{i+1}$ be the pieces that $p_j$ visits from $S_i$ to $S_{i+1}$, in this order, and let $e_m$ be the edge of $p_j$ between $X_m$ and $X_{m+1}$.
The pieces $X_1, \dots, X_r$ are links, and by \cref{lem:pieces}(iii) the edge $e_m$ leaves $X_m$ and enters $X_{m+1}$.
Hence $e_0, X_1, e_1, \dots, X_r, e_r$ is a general path from $S_i$ to $S_{i+1}$ whose edges lie on $p_j$.
It is the only one, because exactly one edge of $p_j$ leaves $S_i$ and a general path is determined by its first edge.
Indeed, each of its links is the piece that contains the head of the preceding edge, and each further edge is the only edge that leaves the preceding link (\cref{lem:degrees}).
Every edge of $F$ lies on $p_1$ or $p_2$ (\cref{lem:orientation}(i)) and every link is visited by $p_1$ or $p_2$ (\cref{lem:degrees}).
Since $p_j$ leads from $S_1$ to $S_k$, each of these edges and links belongs to $P_{i,j}$ for some $i$ and $j$.
An edge of $F$ lies on only one of the two paths and between only one pair of consecutive squares, and a link is entered by only one edge, so $i$ and $j$ are unique.

(iv) Let $R$ be a $u$--$v$ path and let $1 \leq i < k$.
By (ii) and \cref{lem:pieces}(iii), $R$ contains exactly one edge that leaves $S_i$.
It is the first edge of $P_{i,j}$ for some $j$.
Let $P_{i,j} = e_0, L_1, e_1, \dots, L_r, e_r$.
If $R$ contains $e_{m-1}$ for some $m$ with $1 \leq m \leq r$, then $R$ visits $L_m$, and by \cref{lem:pieces}(iii) it contains the only edge that leaves $L_m$, which is $e_m$.
Hence $R$ contains all edges and visits all links of $P_{i,j}$.
Let $e'_0, L'_1, e'_1, \dots$ be the other general path from $S_i$ to $S_{i+1}$.
If $R$ visited a link $L'_m$, it would contain the only edge that enters it, which is $e'_{m-1}$.
If $R$ contained an edge $e'_m$ with $m \geq 1$, it would visit $L'_m$, which contains the tail of $e'_m$.
By induction on $m$, the path $R$ would contain $e'_0$, a second edge that leaves $S_i$.
Hence $R$ contains no edge and visits no link of this general path.

(v) Suppose that $c(s, t) \leq 2$ for two distinct terminals $s, t$ of a square $S_i$.
Let $B$ be the vertex set of the component that contains $s$ in the graph obtained from $G$ by removing a minimum $s$--$t$ cut.
Every edge with exactly one end in $B$ belongs to this cut.
Then $s \in B$, $t \notin B$ and $|\delta(B)| \leq 2$.
If exactly one of $u, v$ belonged to $B$, then $\delta(B)$ would be a minimum $u$--$v$ cut with $s$ and $t$ on different sides, which contradicts \cref{lem:pieces}(i).
After exchanging $s$ with $t$ and $B$ with $V \setminus B$ if necessary, we may thus assume that $u, v \notin B$.
In particular, $s \notin \{u, v\}$, so $s$ is the head of an edge $e$ that enters $S_i$ or the tail of an edge $e$ that leaves $S_i$.
Let $z$ be the other end of $e$.
In the first case, let $\delta(A)$ be the minimum $u$--$v$ cut formed by the edges that enter $S_i$, and put $W = A$.
In the second case, let $\delta(A)$ be the minimum $u$--$v$ cut formed by the edges that leave $S_i$, and put $W = V \setminus A$.
In both cases $|\delta(W)| = 2$, the set $W$ contains $z$ and exactly one of $u, v$, and $W$ is disjoint from $S_i$.
Since $u, v \notin B$, the set $W \cup B$ contains exactly one of $u, v$, so $|\delta(W \cup B)| \geq 2$.
We show that $|\delta(W \cup B)| \leq 2$.
If $z \notin B$, the edge $e$ joins $z \in W \setminus B$ to $s \in B \setminus W$, and \cref{lem:uncrossing} gives $|\delta(W \cap B)| + |\delta(W \cup B)| \leq 2 + 2 - 2$.
If $z \in B$, then $z \in W \cap B$ and $u, v \notin W \cap B$.
The vertex $z$ lies on $p_1$ or $p_2$, and this path contains at least two edges of $\delta(W \cap B)$, so $|\delta(W \cap B)| \geq 2$.
By \cref{lem:uncrossing}, $|\delta(W \cup B)| \leq 2 + 2 - |\delta(W \cap B)| \leq 2$.
In both cases $\delta(W \cup B)$ is a minimum $u$--$v$ cut.
The vertex $s$ lies in $W \cup B$, and the vertex $t$ lies neither in $B$ nor in $W$, because $t \in S_i$.
This contradicts \cref{lem:pieces}(i).
\end{proof}

We call the sequence $S_1, P_{1,1}, P_{1,2}, S_2, \dots, P_{k-1,1}, P_{k-1,2}, S_k$ the \emph{square decomposition} of $G$ with respect to $u$ and $v$ (\cref{fig:square-decomposition}).
The general path $P_{i,1}$ starts with the edge that leaves $S_i$ at $c_i$ and ends with the edge that enters $S_{i+1}$ at $a_{i+1}$, and $P_{i,2}$ leads in the same way from $d_i$ to $b_{i+1}$.
The squares, their order, the links and the two general paths between consecutive squares are determined by $G$, $u$ and $v$, by \cref{def:pieces} and \cref{thm:structure}(ii).
Only the names $P_{i,1}, P_{i,2}$ and $a_i, b_i, c_i, d_i$ depend on the choice of $p_1$ and $p_2$.
Every vertex lies in exactly one square or link, and every edge lies in a square, in a link or in $F$.

\begin{figure}[tb]
    \centering
    \includegraphics{square-decomposition.pdf}
    \caption{The square decomposition: squares (boxes), links (ellipses) and the directed edges of $F$}
    \label{fig:square-decomposition}
\end{figure}

\begin{example}
\label{ex:decomposition}
In the graph $G_a$ of \cref{ex:counterexamples}, let $p_1 = u, a, t, v$ and $p_2 = u, s, b, v$.
The minimum $u$--$v$ cuts are $\{ua, us\}$, $\{ua, sb\}$, $\{at, bv\}$ and $\{tv, bv\}$, so $F$ contains all edges except $ab$.
The squares are $S_1 = \{u\}$, $S_2 = \{a, b\}$ and $S_3 = \{v\}$, and the links are $\{s\}$ and $\{t\}$ (\cref{fig:decomposition-example}).
The terminals of $S_2$ are $a_2 = c_2 = a$ and $b_2 = d_2 = b$.
The general paths are $P_{1,1} = ua$, $P_{1,2} = us, \{s\}, sb$, $P_{2,1} = at, \{t\}, tv$ and $P_{2,2} = bv$.
\end{example}

\begin{figure}[tb]
    \centering
    \includegraphics{decomposition-example.pdf}
    \caption{The square decomposition of $G_a$ with respect to $u$ and $v$}
    \label{fig:decomposition-example}
\end{figure}

\begin{example}
\label{ex:worked}
Let $G_1$ be the graph with the vertices $u, v, 1, \dots, 8$ and the edges
\[
u1,\ u2,\ 12,\ 13,\ 23,\ 36,\ 24,\ 45,\ 57,\ 67,\ 68,\ 7v,\ 8v
\]
(\cref{fig:worked-example}).
We have $c(u, v) = 2$ and $d(u, v) = 5$, and we take $p_1 = u, 1, 3, 6, 8, v$ and $p_2 = u, 2, 4, 5, 7, v$.
The minimum $u$--$v$ cuts are $\{u1, u2\}$, $\{36, 24\}$, $\{36, 45\}$, $\{36, 57\}$, $\{68, 7v\}$ and $\{8v, 7v\}$, so $F$ contains all edges except $12$, $13$, $23$ and $67$.
The squares are $S_1 = \{u\}$, $S_2 = \{1, 2, 3\}$, $S_3 = \{6, 7\}$ and $S_4 = \{v\}$, and the links are $\{4\}$, $\{5\}$ and $\{8\}$.
The terminals are $a_2 = 1$, $b_2 = d_2 = 2$, $c_2 = 3$, $a_3 = c_3 = 6$ and $b_3 = d_3 = 7$.
The general paths $P_{1,1}$, $P_{1,2}$, $P_{2,1}$ and $P_{3,2}$ are the single edges $u1$, $u2$, $36$ and $7v$, while $P_{2,2} = 24, \{4\}, 45, \{5\}, 57$ and $P_{3,1} = 68, \{8\}, 8v$.
Every vertex outside $S_2 \cup S_3$ has degree two, two vertices of the same square are adjacent, and the cut $\{36, 24\}$ separates $S_2$ from $S_3$.
Hence $G_1 \in \DC$.
\end{example}

\begin{figure}[tb]
    \centering
    \includegraphics{worked-example.pdf}
    \caption{The graph $G_1$ of \cref{ex:worked} and its square decomposition}
    \label{fig:worked-example}
\end{figure}

\subsection{Separating Shortest Paths}
\label{sec:separating}

By \cref{lem:basic-bounds} and \cref{prop:easy-pairs}(iii), a pair with $c(u, v) = 2$ satisfies $\cp(u, v) = d(u, v)$ if $G$ contains a separating shortest $u$--$v$ path, and $\cp(u, v) = d(u, v) + 1$ otherwise.
Whether a $u$--$v$ path separates $u$ from $v$ is decided inside the squares.

Let $R$ be a $u$--$v$ path.
By \cref{thm:structure}(ii) and \cref{lem:pieces}(iii), the path $R$ contains exactly one of the two edges that enter $S_i$, for $2 \leq i \leq k$, and exactly one of the two edges that leave $S_i$, for $1 \leq i < k$.
Let $x_i$ be the head of the edge of $R$ that enters $S_i$, and let $\alpha_i$ be the head of the other edge that enters $S_i$.
Let $y_i$ be the tail of the edge of $R$ that leaves $S_i$, and let $\beta_i$ be the tail of the other edge that leaves $S_i$.
We put $x_1 = \alpha_1 = u$ and $y_k = \beta_k = v$.
Then $(x_i, \alpha_i)$ is $(a_i, b_i)$ or $(b_i, a_i)$, the pair $(y_i, \beta_i)$ is $(c_i, d_i)$ or $(d_i, c_i)$, and the trace of $R$ in $S_i$ is an $x_i$--$y_i$ path (\cref{fig:square-terminals}).

\begin{figure}[tb]
    \centering
    \includegraphics{square-terminals.pdf}
    \caption{A square $S_i$ and a trace $x, w, y$ that separates $\alpha$ from $\beta$ in (a) and does not in (b)}
    \label{fig:square-terminals}
\end{figure}

\begin{lemma}[Separation criterion]
\label{lem:criterion}
Let $c(u, v) = 2$, and let $R$ be a $u$--$v$ path.
The graph $G \setminus R$ contains a $u$--$v$ path if and only if, for every $i \in \{1, \dots, k\}$, the vertices $\alpha_i$ and $\beta_i$ lie in the same component of $S_i \setminus R$.
\end{lemma}

\begin{proof}
Let $1 \leq i < k$.
By \cref{thm:structure}(iv), $R$ contains the edges of one of the general paths from $S_i$ to $S_{i+1}$, and it contains no edge and visits no link of the other one, which we denote by $e_0, L_1, e_1, \dots, L_r, e_r$.
The edge $e_0$ is the edge that leaves $S_i$ and does not lie on $R$, so its tail is $\beta_i$, and likewise the head of $e_r$ is $\alpha_{i+1}$.
The links $L_1, \dots, L_r$ are connected and contain no edge of $R$.
Hence $\beta_i$ and $\alpha_{i+1}$ lie in the same component of $G \setminus R$.

Suppose that $\alpha_i$ and $\beta_i$ lie in the same component of $S_i \setminus R$ for every $i$.
Then $u = \alpha_1$, $\beta_1$, $\alpha_2$, $\beta_2$, \dots, $\alpha_k$, $\beta_k = v$ all lie in the same component of $G \setminus R$.

Conversely, let $Q$ be a $u$--$v$ path of $G \setminus R$, and let $1 \leq i \leq k$.
By \cref{thm:structure}(ii), $Q$ visits $S_i$.
If $i \geq 2$, the path $Q$ contains an edge that enters $S_i$ and does not lie on $R$, so the first vertex of $Q$ in $S_i$ is $\alpha_i$.
If $i = 1$, this vertex is $u = \alpha_1$.
In the same way, the last vertex of $Q$ in $S_i$ is $\beta_i$.
The trace of $Q$ in $S_i$ is therefore an $\alpha_i$--$\beta_i$ path of $S_i \setminus R$.
\end{proof}

\begin{theorem}[Separating shortest paths]
\label{thm:characterization}
Let $G$ be a connected graph, let $u$ and $v$ be two vertices of $G$ with $c(u, v) = 2$, and let $S_1, \dots, S_k$ be the squares of $G$ with respect to $u$ and $v$, with the terminals $a_i, b_i, c_i, d_i$.
Then $\cp(u, v) = d(u, v)$ if and only if there are an index $i$, a pair $(x, \alpha) \in \{(a_i, b_i), (b_i, a_i)\}$, a pair $(y, \beta) \in \{(c_i, d_i), (d_i, c_i)\}$ and an $x$--$y$ path $p$ of $S_i$ such that
(i) $d(u, x) + |p| + d(y, v) = d(u, v)$ and
(ii) the vertices $\alpha$ and $\beta$ lie in different components of $S_i \setminus p$.
Otherwise $\cp(u, v) = d(u, v) + 1$.
\end{theorem}

\begin{proof}
Suppose that $i$, $(x, \alpha)$, $(y, \beta)$ and $p$ satisfy (i) and (ii).
Let $R_x$ be a shortest $u$--$x$ path and let $R_y$ be a shortest $y$--$v$ path of $G$.
By (i), the paths $R_x$, $p$ and $R_y$ form a walk from $u$ to $v$ with $d(u, v)$ edges.
Such a walk repeats no vertex, so it is a shortest $u$--$v$ path $R$.
Let $x_i, \alpha_i, y_i, \beta_i$ be defined for $R$ as above, and let $p'$ be the trace of $R$ in $S_i$, which contains $p$.
We show that $\alpha_i$ and $\alpha$ lie in the same component of $S_i \setminus p$.
This holds if $\alpha_i = \alpha$, in particular if $i = 1$.
Otherwise $(x_i, \alpha_i) = (\alpha, x)$, and the part of $p'$ from $x_i = \alpha$ to $x$ is a path of $S_i$ that shares no edge with $p$ and joins $\alpha$ to $\alpha_i = x$.
In the same way, $\beta_i$ and $\beta$ lie in the same component of $S_i \setminus p$.
By (ii), the vertices $\alpha_i$ and $\beta_i$ lie in different components of $S_i \setminus p$, and hence of its subgraph $S_i \setminus R$.
By \cref{lem:criterion}, $R$ is a separating shortest $u$--$v$ path, and $\cp(u, v) = d(u, v)$ by \cref{lem:basic-bounds}.

Conversely, suppose that $\cp(u, v) = d(u, v)$.
By \cref{lem:basic-bounds}, $G$ contains a separating shortest $u$--$v$ path $R$, and by \cref{lem:criterion}, the vertices $\alpha_i$ and $\beta_i$ lie in different components of $S_i \setminus R$ for some $i$.
Let $p$ be the trace of $R$ in $S_i$, and let $(x, \alpha) = (x_i, \alpha_i)$ and $(y, \beta) = (y_i, \beta_i)$.
Condition (ii) holds because $S_i \setminus R = S_i \setminus p$.
Condition (i) holds because the parts of the shortest path $R$ from $u$ to $x$ and from $y$ to $v$ are shortest paths.

The last statement follows from \cref{prop:easy-pairs}(iii).
\end{proof}

A path $p$ with (i) is a shortest $x$--$y$ path of $G$, because $d(u, v) \leq d(u, x) + d(x, y) + d(y, v)$.
A path $p$ with (ii) has at least one edge, because $S_i$ is connected.
The first part of the proof constructs a separating shortest $u$--$v$ path from $p$.

\begin{example}
\label{ex:criterion}
In the graph $G_a$ (\cref{ex:decomposition}), take $i = 2$, $(x, \alpha) = (a, b)$, $(y, \beta) = (b, a)$ and $p = a, b$.
Then $d(u, a) + |p| + d(b, v) = 3 = d(u, v)$, and the graph $S_2 \setminus p$ has no edge, so it separates $\alpha = b$ from $\beta = a$.
Hence $\cp(u, v) = 3$.

In the graph $G_1$ (\cref{ex:worked}), take $i = 3$, $(x, \alpha) = (6, 7)$, $(y, \beta) = (7, 6)$ and $p = 6, 7$.
Then $d(u, 6) + |p| + d(7, v) = 3 + 1 + 1 = d(u, v)$, and $S_3 \setminus p$ has no edge.
Hence $\cp(u, v) = 5$, and $u, 1, 3, 6, 7, v$ is a minimum cut-path.

Let $G_1'$ be obtained from $G_1$ by adding a vertex $9$ adjacent to $6$ and $7$.
The set $F$, the distances from $u$ and to $v$ and the squares $S_1, S_2, S_4$ do not change, and $S_3 = \{6, 7, 9\}$.
The only paths with at least one edge that satisfy (i) are the edges $13$ and $23$ of $S_2$ and the edge $67$ of $S_3$, and each of the graphs $S_2 \setminus \{13\}$, $S_2 \setminus \{23\}$ and $S_3 \setminus \{67\}$ is connected.
Hence $\cp(u, v) = 6$ in $G_1'$.

Let $G_1''$ be obtained from $G_1$ by removing the vertex $5$ and adding the edge $47$.
Then $d(u, v) = 4$, the squares are those of $G_1$, and no path of a square with at least one edge satisfies (i).
Hence $\cp(u, v) = 5$ in $G_1''$.
\end{example}

\subsection{A Linear-Time Algorithm}
\label{sec:algorithm}

In a graph of $\DC$, the paths $p$ of \cref{thm:characterization} have at most two edges, and a constant number of them has to be tested in each square.

\begin{corollary}
\label{cor:two-edges}
Let $G \in \DC$ and let $u, v$ be two vertices of $G$ with $c(u, v) = 2$.
Then every shortest $u$--$v$ path has at most two edges in every square.
Every path $p$ that satisfies condition (i) of \cref{thm:characterization} has $|p| = d(x, y) \leq 2$.
\end{corollary}

\begin{proof}
Let $x$ be an entering terminal and let $y$ be a leaving terminal of a square $S_i$.
If $x \neq y$, then $c(x, y) \geq 3$ by \cref{thm:structure}(v), and $d(x, y) \leq 2$ by \cref{def:class}.
Hence $d(x, y) \leq 2$ in all cases.
The trace of a shortest $u$--$v$ path in $S_i$ is a shortest path of $G$ from an entering terminal to a leaving terminal of $S_i$.
It contains all edges of the path in $S_i$, because an edge of a path joins two consecutive vertices of the path.
A path $p$ with (i) is a shortest $x$--$y$ path of $G$.
\end{proof}

The set $F$ is computed from $p_1$ and $p_2$.
Let $D$ be the directed graph with the vertex set $V$ that contains, for every edge $xy$ of $G$ outside $p_1 \cup p_2$, the arcs from $x$ to $y$ and from $y$ to $x$, and for every edge $xy$ that $p_1$ or $p_2$ traverses from $x$ to $y$, the arc from $y$ to $x$.
It is the residual graph of the flow that sends one unit along $p_1$ and one along $p_2$.
For a directed network, statement (i) of the following lemma is the theorem of Picard and Queyranne~\cite{PicardQueyranne1980} on the closed sets of the residual graph of a maximum flow, and statement (ii) is their criterion for saturated arcs.
We give the short proofs for the flow of value two in the undirected graph $G$, where every edge outside $p_1 \cup p_2$ has both arcs in $D$.

\begin{lemma}
\label{lem:residual}
Let $c(u, v) = 2$.
(i) For a set $A \subseteq V$ with $u \in A$ and $v \notin A$, the set $\delta(A)$ is a minimum $u$--$v$ cut if and only if no arc of $D$ leads from $A$ to $V \setminus A$.
(ii) An edge $xy$ that $p_1$ or $p_2$ traverses from $x$ to $y$ belongs to $F$ if and only if $x$ and $y$ lie in different strong components of $D$.
\end{lemma}

\begin{proof}
(i) Let $\delta(A)$ be a minimum $u$--$v$ cut.
By \cref{lem:orientation}, it consists of an edge of $p_1$ and an edge of $p_2$, and each of them is traversed by its path from $A$ to $V \setminus A$.
The arcs of these two edges lead into $A$, and no other arc has its ends on different sides.
Conversely, suppose that no arc of $D$ leads from $A$ to $V \setminus A$.
Then every edge of $\delta(A)$ lies on $p_1$ or $p_2$ and is traversed by its path from $A$ to $V \setminus A$.
A path that starts in $A$ and has left $A$ cannot leave $A$ again without returning to $A$ first, so each of $p_1$ and $p_2$ contains exactly one edge of $\delta(A)$.
Hence $\delta(A)$ is a $u$--$v$ cut with two edges.

(ii) Let $p_j$ be the path that contains $xy$.
If $xy \in F$, it lies in a minimum $u$--$v$ cut $\delta(A)$, and $x \in A$, $y \notin A$ by \cref{lem:orientation}(ii).
By (i), no directed path of $D$ leads from $x$ to $y$, so $x$ and $y$ lie in different strong components.
Conversely, suppose that $x$ and $y$ lie in different strong components.
Since $D$ contains the arc from $y$ to $x$, no directed path of $D$ leads from $x$ to $y$.
Let $A$ be the set of all vertices that can be reached from $x$ by directed paths of $D$.
Then $x \in A$, $y \notin A$, and no arc leads from $A$ to $V \setminus A$.
The arcs of the edges of $p_j$ form a directed path from $v$ through $y$ and $x$ to $u$.
Hence $u \in A$, and $v \notin A$ because $y$ can be reached from $v$.
By (i), $\delta(A)$ is a minimum $u$--$v$ cut, and it contains $xy$.
\end{proof}

\begin{algorithm}[tb]
\caption{A minimum cut-path if $c(u, v) = 2$, or two vertices that show $G \notin \DC$}
\label{alg:cut-two}
\begin{algorithmic}[1]
\Require a connected graph $G = (V, E)$ and two vertices $u, v$ with $c(u, v) = 2$
\Ensure $\cp(u, v)$ and a minimum cut-path between $u$ and $v$, or two vertices $x, y$ with $c(x, y) \geq 3$ and $d(x, y) \geq 3$ and a cut-path between $u$ and $v$ with $d(u, v) + 1$ edges
\State compute $d(u, z)$ and $d(z, v)$ for all $z \in V$ by two breadth-first searches \label{line:bfs}
\State compute two edge-disjoint $u$--$v$ paths $p_1$ and $p_2$ \label{line:paths}
\State $F \gets$ the edges of $p_1 \cup p_2$ whose ends lie in different strong components of $D$ \label{line:residual}
\State compute the pieces, the squares $S_1, \dots, S_k$ and their terminals \label{line:pieces}
\For{$i \gets 1$ to $k$} \label{line:loop}
    \For{each $(x, \alpha) \in \{(a_i, b_i), (b_i, a_i)\}$ and $(y, \beta) \in \{(c_i, d_i), (d_i, c_i)\}$ with $x \neq y$ and $\alpha \neq \beta$}
        \State $\ell \gets d(u, v) - d(u, x) - d(y, v)$ \label{line:length}
        \If{$\ell \geq 3$}
            \State record the pair $x, y$ \Comment{$c(x, y) \geq 3$ and $d(x, y) \geq 3$} \label{line:record}
        \EndIf
        \State $\mathcal{P} \gets \emptyset$
        \If{$\ell = 1$ and $xy \in E$}
            \State $\mathcal{P} \gets$ the set that contains the path $x, y$ \label{line:edge}
        \ElsIf{$\ell = 2$}
            \State $q \gets$ an $\alpha$--$\beta$ path of $S_i$ \label{line:q}
            \State $\mathcal{P} \gets$ the set of the paths $x, w, y$ of $S_i$ such that $xw$ or $wy$ is an edge of $q$ \label{line:candidates}
        \EndIf
        \For{each $p \in \mathcal{P}$}
            \If{$\alpha$ and $\beta$ lie in different components of $S_i \setminus p$} \label{line:test}
                \State \Return $d(u, v)$ and $R_x \cup p \cup R_y$ for a shortest $u$--$x$ path $R_x$ and a shortest $y$--$v$ path $R_y$ \label{line:found}
            \EndIf
        \EndFor
    \EndFor
\EndFor
\State $R \gets$ a shortest $u$--$v$ path, $C \gets$ the set of the two edges that leave $S_1$ \label{line:union}
\If{a pair $x, y$ is recorded}
    \State \Return $x, y$ and $R \cup C$ \label{line:certificate}
\EndIf
\State \Return $d(u, v) + 1$ and $R \cup C$ \label{line:not-found}
\end{algorithmic}
\end{algorithm}

\Cref{alg:cut-two} computes the square decomposition and searches every square for a path $p$ as in \cref{thm:characterization} with at most two edges.
For a path with two edges it uses that a path which separates $\alpha$ from $\beta$ shares an edge with every $\alpha$--$\beta$ path of the square.
The algorithm does not test whether $G \in \DC$.

\begin{theorem}
\label{thm:linear-time}
There is an algorithm that, given a connected graph $G = (V, E)$ and two distinct vertices $u, v$ of $G$, decides in $O(|V| + |E|)$ time whether $c(u, v) \leq 2$.
If so, it returns within the same time $\cp(u, v)$ and a minimum cut-path between $u$ and $v$, or it returns two vertices $x, y$ with $c(x, y) \geq 3$ and $d(x, y) \geq 3$, which show that $G \notin \DC$, and a cut-path between $u$ and $v$ with $d(u, v) + 1 \leq \cp(u, v) + 1$ edges.
The second outcome does not occur for $G \in \DC$.
\end{theorem}

\begin{proof}
The algorithm first computes, by at most three augmentations of the algorithm of Ford and Fulkerson~\cite[Section~10.4]{Erickson2019Algorithms}, a flow from $u$ to $v$ of value at most three in which every edge of $G$ carries at most one unit.
Each augmentation is a search in the residual graph of the current flow and takes $O(|V| + |E|)$ time.
The first augmentation succeeds because $G$ is connected.
The cut-value $c(u, v)$ is one if the second augmentation fails, two if the third one fails, and at least three otherwise.
If $c(u, v) = 1$, the algorithm returns $d(u, v)$ and a shortest $u$--$v$ path, which is correct by \cref{prop:easy-pairs}(ii).

Let $c(u, v) = 2$.
The two augmenting paths need not be edge-disjoint, because the second one can traverse an edge of the first one backwards, which cancels the flow on this edge.
The paths $p_1$ and $p_2$ are therefore read from the flow of value two.
Let $H$ be the directed graph formed by the edges that carry flow, each directed along the flow.
In $H$, the vertex $u$ has two more leaving than entering arcs, the vertex $v$ has two more entering than leaving arcs, and every other vertex has as many entering as leaving arcs.
No arc of $H$ leaves the set $X$ of the vertices that can be reached from $u$, so the vertices of $X$ have in total at most as many leaving as entering arcs, and hence $v \in X$.
A search in $H$ therefore finds a $u$--$v$ path $p_1$.
After the removal of the arcs of $p_1$, the same argument with one in place of two gives a $u$--$v$ path $p_2$.
The paths $p_1$ and $p_2$ are edge-disjoint paths of $G$, and the arcs of $H$ that remain form cycles and are not used.
Line~\ref{line:paths} of \cref{alg:cut-two} takes these two paths, and the algorithm runs \cref{alg:cut-two}.

\emph{Correctness.}
By \cref{lem:residual}(ii), line~\ref{line:residual} computes the set $F$, and the direction in which $p_1$ or $p_2$ traverses an edge of $F$ gives its tail and its head.
Line~\ref{line:pieces} therefore computes the squares and the terminals of \cref{def:pieces,def:decomposition}.
Every path $p$ that is tested in line~\ref{line:test} is an $x$--$y$ path of $S_i$ with $\ell$ edges, so it satisfies condition (i) of \cref{thm:characterization}.
If the test succeeds, $p$ also satisfies condition (ii).
By the proof of \cref{thm:characterization}, the set returned in line~\ref{line:found} is then a separating shortest $u$--$v$ path, which is a cut-path with $d(u, v) = \cp(u, v)$ edges.
A path returned in line~\ref{line:found} is thus a minimum cut-path in every connected graph.

If line~\ref{line:length} gives $\ell \geq 3$ for two terminals $x \neq y$, then $d(x, y) \geq \ell \geq 3$, because $d(u, v) \leq d(u, x) + d(x, y) + d(y, v)$, and $c(x, y) \geq 3$ by \cref{thm:structure}(v).
A pair recorded in line~\ref{line:record} therefore shows that $G \notin \DC$, and no pair is recorded if $G \in \DC$.

Suppose now that $\cp(u, v) = d(u, v)$ and that no pair is recorded.
By \cref{thm:characterization}, there are $i$, $(x, \alpha)$, $(y, \beta)$ and $p$ with (i) and (ii).
Then $\alpha \neq \beta$ and the path $p$ has at least one edge, so $x \neq y$, and line~\ref{line:length} gives $\ell = |p|$ for these terminals.
Since the pair $x, y$ is not recorded, $\ell \in \{1, 2\}$.
If $\ell = 1$, then $p$ is the path of line~\ref{line:edge}.
If $\ell = 2$, let $p = x, w, y$.
The path $q$ of line~\ref{line:q} is not a path of $S_i \setminus p$, so it contains the edge $xw$ or the edge $wy$.
In both cases $p \in \mathcal{P}$, so the algorithm returns in line~\ref{line:found}, for this path or for an earlier one.

Hence, if the algorithm reaches line~\ref{line:not-found}, then $\cp(u, v) = d(u, v) + 1$.
The set $C$ of line~\ref{line:union} is a minimum $u$--$v$ cut by \cref{thm:structure}(i), and $R$ contains exactly one of its edges by \cref{lem:orientation}(ii).
Hence $R \cup C$ is a cut-path with $d(u, v) + 1$ edges, which is a minimum cut-path in line~\ref{line:not-found} and has at most $\cp(u, v) + 1$ edges in line~\ref{line:certificate}.

\emph{Running time.}
The graph is given by adjacency lists.
Breadth-first search~\cite[Section~8.4]{Erickson2019Algorithms} and the computation of strong components~\cite{Tarjan1972DFS} take $O(|V| + |E|)$ time.
The paths $p_1$ and $p_2$ are found by two searches in $H$.
Lines~\ref{line:bfs} to~\ref{line:pieces} consist of a constant number of such computations and of one search in $G \setminus F$ for the pieces.
Let $S_i$ have $n_i$ vertices and $m_i$ edges.
For each of the at most four choices of $(x, \alpha)$ and $(y, \beta)$, the path $q$ is found by one search in $S_i$.
The path $q$ has at most two edges at $x$ and at most two edges at $y$, so $\mathcal{P}$ has at most four members.
Whether $x, w, y$ is a path of $S_i$ is decided for each of these vertices $w$ after marking the neighbors of $x$ and of $y$.
A vertex of $S_i$ is an end of at most four edges outside $S_i$, the edges that enter or leave $S_i$, so its neighbors are marked in $O(m_i + 1)$ time.
The test of line~\ref{line:test} is one search in $S_i \setminus p$.
The square $S_i$ therefore takes $O(n_i + m_i)$ time.
The squares are pairwise disjoint, so the loop of line~\ref{line:loop} takes $O(|V| + |E|)$ time.
The paths $R_x$, $R_y$ and $R$ are read from the breadth-first search trees.
\end{proof}

\Cref{alg:cut-two} returns $\cp(u, v)$ and a minimum cut-path whenever no pair is recorded, in particular if every entering terminal of every square is at distance at most two from every leaving terminal of the same square.

\begin{example}
\label{ex:two-edges}
Let $G_a^{+}$ be obtained from the graph $G_a$ of \cref{ex:counterexamples} by adding a vertex $r$ adjacent to $u$ and $b$.
The vertices of degree greater than two are $u$, $a$ and $b$, which are pairwise at distance at most two, so $G_a^{+} \in \DC$.
The minimum $u$--$v$ cuts are $\{at, bv\}$ and $\{tv, bv\}$.
Hence the squares are $S_1 = \{u, s, r, a, b\}$ and $S_2 = \{v\}$, the only link is $\{t\}$, and $c_1 = a$ and $d_1 = b$ for $p_1 = u, a, t, v$.
For $(x, \alpha) = (u, u)$ and $(y, \beta) = (b, a)$, \cref{alg:cut-two} computes $\ell = 3 - 0 - 1 = 2$.
Let $q$ be the path $u, a$.
Its only edge is $ua$, so $w = a$, and $\mathcal{P}$ contains the path $p = u, a, b$.
The vertex $a$ is isolated in $S_1 \setminus p$, so the algorithm returns $3$ and the path $u, a, b, v$.
\end{example}

\begin{corollary}
\label{cor:class-pairs}
Let $G \in \DC$ and let $u, v$ be two distinct vertices of $G$.
If $c(u, v) \leq 2$, in particular if $d(u, v) \geq 3$, then $\cp(u, v)$ and a minimum cut-path between $u$ and $v$ are computed in $O(|V| + |E|)$ time.
If $d(u, v) = 1$, they are computed by one computation of a minimum $u$--$v$ cut.
\end{corollary}

\begin{proof}
A pair with $d(u, v) \geq 3$ has $c(u, v) \leq 2$ by \cref{def:class}.
The first statement is \cref{thm:linear-time}, whose second outcome does not occur for $G \in \DC$, and the second is \cref{prop:easy-pairs}(i).
A minimum $u$--$v$ cut is computed in polynomial time by a maximum flow~\cite[Sections~10.3 and~10.4]{Erickson2019Algorithms}.
\end{proof}

\begin{remark}[Coarser decompositions]
\label{rem:coarser}
Let $F'$ be the union of a nonempty set of minimum $u$--$v$ cuts, and define tails, heads, pieces, squares, links, terminals and general paths with $F'$ in place of $F$.
\Cref{lem:orientation,lem:pieces,lem:degrees}, statements (i) to (iv) of \cref{thm:structure}, \cref{lem:criterion} and \cref{thm:characterization} remain true if \cref{lem:pieces}(i) is restricted to the minimum $u$--$v$ cuts contained in $F'$.
Their proofs use only that every edge of $F$ lies in a minimum $u$--$v$ cut contained in $F$.
Statement (v) of \cref{thm:structure} uses all minimum $u$--$v$ cuts and can fail for $F'$.
If every entering terminal of every square of the coarser decomposition is at distance at most two from every leaving terminal of the same square, then lines~\ref{line:loop} to~\ref{line:not-found} of \cref{alg:cut-two} are correct for this decomposition as well.
The square decomposition defined here is the finest of the decompositions of this remark.
For $G \in \DC$, it satisfies the distance condition by \cref{thm:structure}(v).
\end{remark}

